\chapter[1947 --- Cori et al.]{1947: Carl Ferdinand Cori, Gerty Theresa Cori, and Bernardo Houssay}
\label{ch:1947_Cori}

\section*{Main Contribution}

\textit{Elucidated the metabolic pathways by which the body converts glycogen to glucose, explaining how the body maintains blood sugar levels.}\cite{nobel-1947,lecture-1947}

\section{Official Citation}

for their discovery of the course of the catalytic conversion of glycogen; and for his discovery of the part played by the hormone of the anterior pituitary lobe in the metabolism of sugar

\section{Background and Historical Context}

\subsection{Carl Ferdinand Cori}

Carl Ferdinand Cori was born on February 5, 1896, in Prague, Austria-Hungary (now the Czech Republic), into a distinguished family with a strong tradition of scientific and intellectual achievement. His father was a biologist, and young Carl grew up in an environment that valued scientific inquiry. He attended the German University of Prague, where he studied medicine and developed an early interest in biochemistry. During World War I, he served as a medical officer in the Austrian Army, an experience that deepened his interest in the physiological mechanisms underlying disease. After the war, he returned to Prague and began his research career at the University of Prague's Institute of Pharmacology.

\subsection{Gerty Theresa Cori}

Gerty Theresa Cori was born on August 15, 1896, in Prague, Austria-Hungary. She was one of the first women to attend medical school in the Austro-Hungarian Empire and graduated with her medical degree from the German University of Prague in 1920. She married Carl Cori in 1920, and the couple would become one of the most notable husband-and-wife teams in the history of science. Gerty's expertise in enzymology and metabolic biochemistry complemented Carl's strengths in physiology, and together they would make fundamental contributions to our understanding of carbohydrate metabolism.

\subsection{Bernardo Houssay}

Bernardo Houssay was born on April 10, 1887, in Buenos Aires, Argentina, into a family of French descent. He studied medicine at the University of Buenos Aires, where he graduated in 1911, and subsequently joined the university's Department of Physiology as an assistant. Houssay developed a deep interest in the endocrine system---the network of glands that produce hormones---and he would become one of the world's leading authorities on the hormonal regulation of metabolism. His research was conducted entirely in Argentina, making him the first Latin American scientist to receive a Nobel Prize in Physiology or Medicine.

The scientific context of the 1930s and 1940s was one of rapid progress in biochemistry and endocrinology. The discovery of insulin by Banting and Best in 1921 had transformed the treatment of diabetes, but the mechanisms by which insulin and other hormones regulate blood sugar levels remained poorly understood. It was within this dynamic intellectual environment that the Coris and Houssay made their landmark discoveries.

\section{The Discovery/Contribution}

The Nobel Prize was awarded to Carl and Gerty Cori jointly ``for their discovery of the course of the catalytic conversion of glycogen,'' and to Bernardo Houssay ``for his discovery of the part played by the hormone of the anterior pituitary lobe in the metabolism of sugar.'' Together, their work elucidated the fundamental mechanisms by which the body stores, mobilizes, and utilizes glucose---the primary fuel of human cells.

\subsection{The Cori Cycle}

Carl and Gerty Cori's primary contribution was the elucidation of the metabolic pathway that bears their name: the Cori cycle (also known as the lactic acid cycle). This cycle describes the interconversion of glucose and lactate (lactic acid) between the liver and skeletal muscle.

In the Cori cycle, glucose is taken up by muscle cells and broken down through glycolysis to produce pyruvate, which is then converted to lactate. The lactate is released into the bloodstream and transported to the liver, where it is converted back to glucose through gluconeogenesis. The glucose is then released into the bloodstream and transported back to the muscles, where it can be used as fuel for contraction. This cycle allows the body to regenerate glucose from the lactate produced during anaerobic metabolism, ensuring a continuous supply of fuel to the muscles during exercise.

The Coris's elucidation of this cycle was a tour de force of experimental biochemistry. Working at the State Institute for the Study of Malignant Disease in Buffalo, New York (now the Roswell Park Comprehensive Cancer Center), they used a combination of physiological experiments and chemical analyses to trace the fate of glucose through the body. Their work demonstrated that the conversion of glucose to lactate and the reverse reaction of lactate to glucose are catalyzed by different enzymes in different tissues---a finding that explained how the body can maintain the cycle in a continuous state of flux.

\subsection{Glycogen Phosphorylase and the Mechanism of Glycogen Breakdown}

The Coris also made a fundamental contribution to our understanding of glycogen metabolism with the discovery of the enzyme glycogen phosphorylase. Glycogen is a branched polymer of glucose that serves as the body's primary short-term energy reserve, stored in the liver and skeletal muscles. When energy is needed, glycogen is broken down to release glucose molecules.

The Coris demonstrated that this breakdown is catalyzed by glycogen phosphorylase, which cleaves glucose units from the ends of glycogen chains by adding inorganic phosphate (a process called phosphorolysis). The product of this reaction is glucose-1-phosphate, which can be converted to glucose-6-phosphate and then to free glucose for use as fuel.

The discovery of glycogen phosphorylase was significant not only for its immediate contribution to our understanding of carbohydrate metabolism but also for its broader implications for enzymology. The Coris showed that glycogen phosphorylase exists in two forms---an active form (phosphorylase a) and an inactive form (phosphorylase b)---and that the interconversion between these forms is regulated by phosphorylation and dephosphorylation. This was one of the first demonstrations that enzyme activity can be regulated by covalent modification, a principle that has since been found to be fundamental to the regulation of virtually all metabolic pathways.

\subsection{Bernardo Houssay and the Anterior Pituitary Hormone}

Bernardo Houssay's Nobel Prize was awarded for his discovery of the role of the anterior pituitary gland in sugar metabolism. The anterior pituitary---a small endocrine gland located at the base of the brain---produces a variety of hormones that regulate growth, reproduction, and metabolism. Houssay demonstrated that one of these hormones, which he called the ``diabetogenic hormone'' (now known to include growth hormone and adrenocorticotropic hormone), has significant effects on blood sugar levels.

Working with dogs in his laboratory in Buenos Aires, Houssay showed that removal of the anterior pituitary gland (hypophysectomy) in diabetic animals reduced the severity of the diabetes, while the administration of anterior pituitary extracts increased blood sugar levels and exacerbated the diabetic condition. These experiments provided the first direct evidence that the anterior pituitary gland plays a crucial role in the regulation of blood sugar metabolism.

Houssay's work was particularly significant because it demonstrated that diabetes is not solely a disease of the pancreas (as had been assumed since the discovery of insulin) but involves a complex interplay between multiple endocrine glands. The anterior pituitary hormone acts, at least in part, by opposing the effects of insulin and by stimulating the breakdown of glycogen in the liver. This insight laid the groundwork for our modern understanding of the hormonal regulation of blood sugar levels and the pathophysiology of diabetes.

\section{Impact on Modern Medicine}

The discoveries of the Coris and Houssay have had significant implications for modern medicine, particularly in the fields of endocrinology, metabolic disease, and sports medicine.

\subsection{Understanding and Treating Diabetes}

The Coris's elucidation of the Cori cycle and glycogen metabolism, combined with Houssay's discovery of the anterior pituitary's role in sugar metabolism, fundamentally changed our understanding of diabetes. Diabetes is characterized by the inability of the body to properly regulate blood sugar levels, leading to chronic hyperglycemia (high blood sugar). The work of the Coris and Houssay revealed that this dysregulation involves not only the pancreas and its production of insulin but also the liver, skeletal muscle, and multiple endocrine glands.

This understanding has been essential for the development of modern diabetes management. The recognition that the liver plays a central role in glucose production has led to the development of drugs such as metformin, which reduces hepatic glucose output and is the first-line treatment for type 2 diabetes. Understanding the hormonal regulation of blood sugar has also been crucial for the development of insulin analogs and incretin-based therapies that are used to treat diabetes today.

\subsection{Glycogen Storage Diseases}

The Coris's work on glycogen metabolism has been essential for the diagnosis and treatment of glycogen storage diseases (GSDs), a group of inherited metabolic disorders caused by defects in the enzymes involved in glycogen synthesis or breakdown. These diseases, which include Von Gierke disease (GSD type I), Pompe disease (GSD type II), and McArdle disease (GSD type V), cause a range of symptoms including muscle weakness, liver enlargement, and hypoglycemia.

Understanding the biochemistry of glycogen metabolism, as elucidated by the Coris, has been essential for the diagnosis of these diseases and for the development of treatment strategies. Enzyme replacement therapy, in which the deficient enzyme is provided exogenously, has been developed for some GSDs (notably Pompe disease) and represents a direct application of the Coris's basic research to clinical medicine.

\subsection{Exercise Physiology and Sports Medicine}

The Cori cycle is fundamental to our understanding of exercise physiology. During intense exercise, when the demand for energy exceeds the capacity of aerobic metabolism, muscle cells rely on anaerobic glycolysis, producing large quantities of lactate. The Cori cycle allows this lactate to be converted back to glucose in the liver, sustaining the body's energy supply during prolonged physical activity. Understanding the Cori cycle has been essential for the development of training programs for athletes and for the management of metabolic conditions that affect exercise capacity.

\subsection{Endocrinology and Hormone Therapy}

Houssay's discovery of the anterior pituitary's role in sugar metabolism has had broad implications for endocrinology. The anterior pituitary produces multiple hormones that regulate growth, reproduction, and metabolism, and understanding how these hormones interact has been essential for the diagnosis and treatment of endocrine disorders. Hormone replacement therapy for conditions such as growth hormone deficiency, adrenal insufficiency, and thyroid disease is based on principles that trace back to Houssay's foundational research.

\section{Legacy}

The legacies of Carl and Gerty Cori and Bernardo Houssay are intertwined with the broader history of biochemistry and endocrinology in the twentieth century. The Coris, who were a remarkable husband-and-wife team, demonstrated the power of collaborative research and the importance of persistence in the face of obstacles. After emigrating to the United States in 1922, they faced discrimination---Gerty was not given a faculty position at Washington University in St.\ Louis despite her essential contributions to their joint research---yet they continued to work together with notable productivity and success. Gerty Cori died on October 26, 1957, at the age of 61, from a bone marrow disease that may have been related to radiation exposure during her research.

Carl Cori continued his research after Gerty's death, focusing on the enzymology of carbohydrate metabolism and the biochemistry of cancer. He received numerous honors, including the National Medal of Science, and died on October 20, 1984, at the age of 88.

Bernardo Houssay continued to lead the Department of Physiology at the University of Buenos Aires and became a powerful advocate for science education and research in Argentina. He founded the Argentine National Research Council and played a key role in building Argentina's scientific infrastructure. He died on September 21, 1971, at the age of 84.

The work of the Coris and Houssay has had an enduring impact on our understanding of metabolic disease and endocrine regulation. Their discoveries continue to inform the diagnosis and treatment of diabetes, glycogen storage diseases, and other metabolic conditions, and they remain essential reading for students of biochemistry and physiology around the world.


\section{Modern Developments}

The Cori cycle work has evolved into modern metabolic medicine. Understanding of glycogen metabolism has led to treatments for glycogen storage diseases (enzyme replacement therapy). Lactate, once considered merely a waste product, is now recognized as an important metabolic fuel and signaling molecule. Targeting glycolysis in cancer cells (the Warburg effect) is an active area of drug development. Understanding of glucose metabolism has led to continuous glucose monitoring and closed-loop insulin delivery systems for diabetes.


\section{Bibliography}

\begin{thebibliography}{9}
\bibitem{nobel-1947}
Nobel Prize Outreach, ``The Nobel Prize in Physiology or Medicine 1947,''\emph{NobelPrize.org}.\\
\url{https://www.nobelprize.org/prizes/medicine/1947/summary/}

\bibitem{lecture-1947}
Carl Ferdinand Cori, ``Nobel Lecture,''\emph{NobelPrize.org}. Winner-authored primary source; downloadable from the official Nobel Prize archive.\\
\url{https://www.nobelprize.org/prizes/medicine/1947/cori-cf/lecture/}
\end{thebibliography}
